% ============================================================
%  EXACT DISCRETISATION AND BOUNDARY OBSERVABLES
%  IN LORENTZIAN CAUSAL DIAMONDS  —  v3
%
%  Revised following peer review. Changes from v1:
%  (R1) CD redefined as 2-complex; partial-CD as boundary links
%  (R2) Boundary sum theorem renamed; Stokes label corrected
%  (R3) beta/13 labelled [ANSATZ] with justification
%  (R4) "exact within truncation" language corrected
%  (R5) Finite-system caveat for crossover added
%  (R6) All 21 plaquettes listed in Appendix D
%  (R7) K_bdy written as explicit block matrix
%  (R8) Flat-connection theorem in Section 5
%  (R9) CDT comparison remark added
%  (R10) A_mu constant assumption noted in Theorem 4 proof
%  (R11) "holographic" justified precisely
%  Changes from v2:
%  (R12) Betti-number claim removed; replaced by exact
%        flat-connection count with gauge-theoretic interpretation
%  (R13) CW-complex obstruction explained (star-graph 1-skeleton
%        has no 1-cycles, so plaquettes cannot be attached as
%        topological 2-cells; Euler characteristic cross-check
%        13 - 12 + 21 = 22 ≠ 1 - 4 + 13 = 10 confirms this)
%  (R14) Extended K_{6,6} complex noted for reference
%  (R15) Abstract, introduction, keywords, discussion updated
%  (R16) Appendix D plaquettes 20–21 corrected
%  (R17) K_bdy sign corrected and K_total spectrum updated
% ============================================================

\documentclass[12pt, a4paper]{article}

\usepackage[margin=2.5cm]{geometry}
\usepackage{amsmath, amssymb, amsthm, amsfonts}
\usepackage{mathtools}
\usepackage{bm}
\usepackage{xcolor}
\usepackage[colorlinks=true,linkcolor=black,citecolor=black,urlcolor=black]{hyperref}
\usepackage{booktabs}
\usepackage{array}
\usepackage{longtable}
\usepackage{enumitem}
\usepackage{microtype}
\usepackage{cleveref}
\usepackage{thmtools}
\usepackage{thm-restate}
\usepackage{mathrsfs}

% Theorem environments
\theoremstyle{plain}
\newtheorem{theorem}{Theorem}[section]
\newtheorem{proposition}[theorem]{Proposition}
\newtheorem{lemma}[theorem]{Lemma}
\newtheorem{corollary}[theorem]{Corollary}

\theoremstyle{definition}
\newtheorem{definition}[theorem]{Definition}
\newtheorem{construction}[theorem]{Construction}
\newtheorem{ansatz}[theorem]{Ansatz}

\theoremstyle{remark}
\newtheorem{remark}[theorem]{Remark}
\newtheorem{notation}[theorem]{Notation}

% Macros
\newcommand{\Lat}{\mathscr{L}}
\newcommand{\CD}{\mathscr{D}}
\newcommand{\BdyCD}{\partial\mathscr{D}}
\newcommand{\LL}{\mathcal{L}}
\newcommand{\FS}{\mathcal{F}}
\newcommand{\PS}{\mathcal{P}}
\newcommand{\Plaq}{\mathcal{Q}}
\newcommand{\Nloc}{N_{\mathrm{local}}}
\newcommand{\ssp}{s^2}
\newcommand{\Cstar}{\beta_c}
\newcommand{\Kmat}{\mathbf{K}}
\newcommand{\Mmat}{\mathbf{M}}
\newcommand{\neff}{n^\mu_{\mathrm{eff}}}
\newcommand{\R}{\mathbb{R}}
\newcommand{\Z}{\mathbb{Z}}
\newcommand{\Tr}{\operatorname{Tr}}
\newcommand{\rank}{\operatorname{rank}}
\newcommand{\ZBF}{Z_{\mathrm{BF}}}

\definecolor{ansatzblue}{rgb}{0.0,0.25,0.55}
\newcommand{\ANSATZ}{\textcolor{ansatzblue}{\textbf{[ANSATZ]}}}

\title{\Large\bfseries Exact Discretisation and Boundary Observables\\[4pt]
       in Lorentzian Causal Diamonds}
\author{Yannick Schmitt \\ 
\texttt{ynnk.schmitt@gmail.com} \\
\vspace{4pt} 
\href{https://doi.org/10.5281/zenodo.19338306}{DOI: 10.5281/zenodo.19338306}}
\date{\today}

\begin{document}
\maketitle

\begin{abstract}
We study the ternary Minkowski lattice
$\Lat = \{-1,0,+1\}^4$ with
$\eta = \operatorname{diag}(-1,+1,+1,+1)$
and prove a sequence of exact geometric and algebraic results.
\textbf{Lemma~\ref{thm:12lightlike}} establishes by exhaustive
enumeration that $\Lat$ contains exactly twelve lightlike
nearest-neighbour vectors forming two null sheets of six mutually
spacelike channels each.
\textbf{Proposition~\ref{prop:d4}} identifies these twelve vectors
with the mixed roots of the $D_4$ root system, yielding the exact
Minkowski partition $24 = 12 + 12$.
\textbf{Theorem~\ref{thm:diamond}} proves that the causal diamond
two-complex $\CD$ satisfies five independent geometric conditions;
its boundary $\BdyCD$ is the twelve-link null set, which spans $\R^4$.
\textbf{Theorem~\ref{thm:boundary}} evaluates the discrete boundary
sum of a $U(1)$ gauge field over $\BdyCD$ using the Lorentzian
oriented chain complex, yielding an effective coupling vector
$\neff = (12,0,0,0)$ --- purely temporal.
This is a discrete boundary sum identity; no bulk curvature integral
is computed.
The $U(1)$ lattice BF action on the 21 order-four plaquettes of $\CD$
is constructed with an ansatz boundary coupling at level $\beta/13$.
A leading-order approximate partition function exhibits a
finite-system crossover at $\Cstar \approx 2.7364$.
The plaquette Laplacian has exact integer spectrum
$\{0^{(4)}, 6^{(2)}, 8^{(3)}, 10^{(2)}, 28^{(1)}\}$;
its four-dimensional null space is identified with the space of
flat connections --- link-variable assignments with vanishing
holonomy on every plaquette.
All 21 plaquettes and the $12\times12$ Laplacian are given
explicitly in the appendices.
\end{abstract}

\tableofcontents

\bigskip\noindent
\textbf{MSC 2020:} 83C27, 81T25, 52B11, 81Q35.\\[2pt]
\textbf{Keywords:} ternary Minkowski lattice; causal diamond;
$D_4$ root system; oriented chain complex; $U(1)$ BF theory;
plaquette Laplacian; flat connections.

%------------------------------------------------------------------
\section{Introduction}\label{sec:intro}
%------------------------------------------------------------------

Discrete models of spacetime have motivated substantial literature
in mathematical physics, including causal set
theory~\cite{Bombelli1987,Henson2006}, Regge
calculus~\cite{Regge1961}, causal dynamical
triangulations~\cite{Ambjorn2004}, and spin foam
models~\cite{Perez2013}.
The present paper investigates a concrete, fully explicit object:
the ternary Minkowski lattice $\Lat = \{-1,0,+1\}^4$ with
$\eta = \operatorname{diag}(-1,+1,+1,+1)$.
This is the smallest four-dimensional lattice closed under the
two-valued sign map that contains the $D_4$ root system.

All results are confined to this finite object.
The lattice has $81$ points; all relevant structures are enumerated
exhaustively; every claimed spectrum is the exact integer eigenvalue
set of an explicit matrix.
Five main results are established:
\begin{enumerate}[label=\textbf{(\roman*)}]
  \item $(\Lat,\eta)$ contains exactly twelve lightlike
    nearest-neighbour vectors forming two null sheets
    (\cref{thm:12lightlike}).
  \item These twelve vectors are the mixed roots of $D_4$, and
    $\eta$ partitions all $24$ $D_4$ roots as $12 + 12$
    (\cref{prop:d4}).
  \item The twelve vectors form the boundary $\BdyCD$ of a
    two-complex $\CD$ characterised by five geometric conditions
    (\cref{thm:diamond}).
  \item With the Lorentzian time-orientation on the boundary chain
    complex, the discrete boundary sum of a gauge field over
    $\BdyCD$ evaluates to $\neff = (12,0,0,0)$
    (\cref{thm:boundary}).
  \item The $U(1)$ lattice BF theory on the 21 plaquettes of $\CD$
    has an approximate partition function with finite-system
    crossover $\Cstar \approx 2.7364$, an exact plaquette
    Laplacian with integer spectrum $\{0^4, 6^2, 8^3, 10^2, 28\}$,
    and a four-dimensional space of flat connections
    (\cref{sec:BF}).
\end{enumerate}

No cosmological, phenomenological, or continuum field-theory claims
are made.
The crossover $\Cstar$ is a feature of a 12-link finite system, not
a thermodynamic phase transition.
The term \emph{discrete causal diamond} is used for a finite
combinatorial object satisfying Definition~\ref{def:diamond};
\emph{boundary-spanning} refers to the $\R^4$ spanning condition
(condition~(d) of that definition), not to a dynamical
bulk-boundary duality.

The paper is organised as follows.
\Cref{sec:lattice} proves results (i)--(ii).
\Cref{sec:diamond} defines $\CD$ and proves (iii).
\Cref{sec:boundary} establishes (iv).
\Cref{sec:plaquettes} enumerates the plaquette structure and
computes the flat-connection space.
\Cref{sec:BF} constructs the $U(1)$ BF theory and derives (v).
\Cref{sec:discussion} lists open problems.

%------------------------------------------------------------------
\section{The Ternary Minkowski Lattice}\label{sec:lattice}
%------------------------------------------------------------------

\begin{definition}[Ternary Minkowski Lattice]\label{def:lattice}
The \emph{ternary Minkowski lattice} is
$\Lat = \{-1,0,+1\}^4 \subset \R^4$
with Minkowski inner product
$\eta(u,v) = -u_0 v_0 + u_1 v_1 + u_2 v_2 + u_3 v_3$,
where $u_0$ is the time component and lattice spacing $a=1$.
The Minkowski norm squared of $v \in \Lat$ is $\ssp(v) = \eta(v,v)$.
\end{definition}

\begin{notation}
Write $\Lat^* = \Lat\setminus\{0\}$ for the $80$ non-origin
lattice points, classified as
\emph{lightlike} ($\ssp = 0$),
\emph{spacelike} ($\ssp > 0$), or
\emph{timelike} ($\ssp < 0$).
\end{notation}

\begin{lemma}[Twelve Lightlike Nearest Neighbours]\label{thm:12lightlike}
The lattice $(\Lat,\eta)$ contains exactly twelve lightlike
non-origin vectors:
\[
  \LL = \bigl\{(\pm1,\,\sigma_j\bm{e}_j)
    : j\in\{1,2,3\},\;\sigma_j\in\{-1,+1\}\bigr\}.
\]
The partition of $\Lat^*$ is
$80 = 12_{\ssp=0} + 66_{\ssp>0} + 2_{\ssp<0}$.
\end{lemma}

\begin{proof}
The condition $\ssp(v) = -t^2 + x_1^2+x_2^2+x_3^2 = 0$
with $t, x_j \in \{-1,0,1\}$ and $v \neq 0$
requires $t \neq 0$ (else all $x_j = 0$), so $t^2 = 1$,
forcing exactly one $x_j = \pm 1$:
$2 \times 3 \times 2 = 12$ solutions.
The two purely timelike vectors $(\pm1,0,0,0)$ give $\ssp = -1$;
the remaining $80 - 12 - 2 = 66$ are spacelike.
All 80 vectors are listed in Appendix~\ref{app:enum}.
\qed
\end{proof}

\subsection{Connection to the \texorpdfstring{$D_4$}{D4} Root System}

\begin{definition}[$D_4$ Root System]\label{def:d4}
The $D_4$ root system is
$\Phi(D_4) = \{v\in\Z^4 : |v|^2_\mathrm{Eucl}=2\}$,
consisting of all permutations of $(\pm1,\pm1,0,0)$;
$|\Phi(D_4)| = 24$.
\end{definition}

\begin{proposition}[$D_4$ Root Splitting]\label{prop:d4}
The Minkowski metric $\eta$ partitions $\Phi(D_4)$ as
$24 = 12_\mathrm{lightlike} + 12_\mathrm{spacelike}$.
The 12 lightlike $D_4$ roots coincide exactly with $\LL$.
\end{proposition}

\begin{proof}
A root with $v_0=\pm1$ and exactly one spatial $v_j=\pm1$
satisfies $\ssp=-1+1=0$ (lightlike): $2\times3\times2=12$.
A root with $v_0=0$ and two spatial components satisfies
$\ssp=2>0$ (spacelike): $\binom{3}{2}\cdot4=12$.
No $D_4$ root is timelike under $\eta$.
\qed
\end{proof}

\begin{remark}[$B_4$ vs.\ $\mathrm{SO}(3,1)$]\label{rem:b4symmetry}
The automorphism group of $\Lat$ as a pointed set is the
hyperoctahedral group $B_4$ (order $384$): all permutations
and sign flips of coordinates.
The Lorentz group $\mathrm{SO}(3,1)$ acts on the ambient
Minkowski space $\R^4$, not on the finite lattice points.
Whether $B_4$ invariance implies infrared $\mathrm{SO}(3,1)$
emergence is Open Problem~O4 below.
\end{remark}

%------------------------------------------------------------------
\section{The Discrete Causal Diamond}\label{sec:diamond}
%------------------------------------------------------------------

We define the causal diamond as a \emph{two-complex}
$\CD$ whose boundary $\BdyCD$ is the set of lightlike links.
This distinction between the complex $\CD$ and its boundary
$\BdyCD$ is essential for the boundary integration in \cref{sec:boundary}.

\begin{definition}[Discrete Causal Diamond 2-Complex]\label{def:diamond}
A \emph{discrete causal diamond} in $(\Lat,\eta)$ is a finite
two-complex $\CD$ whose:
\begin{itemize}[noitemsep]
  \item \emph{0-cells}: the origin $\bm{0}$ and the twelve points of $\LL$;
  \item \emph{1-cells}: the twelve null links $\{[\bm{0},L]:L\in\LL\}$;
  \item \emph{2-cells}: the twenty-one order-four closed loops
    (see \cref{sec:plaquettes}).
\end{itemize}
The \emph{boundary} $\BdyCD$ decomposes as
$\BdyCD = \FS \cup \PS$, where
\begin{align*}
  \FS &= \{L\in\LL : L_0=+1\}, \quad |\FS|=6,\\
  \PS &= \{L\in\LL : L_0=-1\}, \quad |\PS|=6.
\end{align*}
The complex $\CD$ is a \emph{boundary-spanning causal diamond}
when $\BdyCD$ satisfies:
\begin{enumerate}[label=(\alph*), noitemsep]
  \item\label{def:a} All boundary links are lightlike.
  \item\label{def:b} Any two distinct links within $\FS$ (or
    both in $\PS$) are spacelike-separated.
  \item\label{def:c} $\FS\cap\PS = \emptyset$.
  \item\label{def:d} $\FS\cup\PS$ spans $\R^4$.
  \item\label{def:e} For each $L_f\in\FS$ there is a unique
    $-L_f\in\PS$ (antipodal pairing).
\end{enumerate}
The qualifier \emph{boundary-spanning} refers to condition~\ref{def:d}:
all directional information about $\R^4$ is encoded on the
twelve boundary null links, in the sense that any vector in
$\R^4$ is a linear combination of elements of $\BdyCD$.
No dynamical bulk-boundary duality of the AdS/CFT type is implied.
\end{definition}

\begin{theorem}[Causal Diamond Boundary]\label{thm:diamond}
The two-complex $\CD$ built on $\LL$ is a boundary-spanning causal
diamond: its boundary $\BdyCD = \FS\cup\PS$ satisfies all
five conditions~\ref{def:a}--\ref{def:e}.
\end{theorem}

\begin{proof}
Condition~\ref{def:a} follows from \cref{thm:12lightlike}.
Condition~\ref{def:b}: for distinct $u,v\in\FS$,
$(u-v)_0=0$ and at least one spatial component differs, so
$\ssp(u-v) = \sum_j(u_j-v_j)^2 \geq 2 > 0$;
all $\binom{6}{2}=15$ pairs are verified spacelike
in Appendix~\ref{app:pairwise}.
Condition~\ref{def:c}: by definition.
Condition~\ref{def:d}: four linearly independent vectors in $\FS$
are exhibited in Appendix~\ref{app:span}.
Condition~\ref{def:e}: the map $L_f\mapsto -L_f$ is a bijection
$\FS\to\PS$.
\qed
\end{proof}

\begin{remark}[Causal Channel Count]\label{rem:bandwidth}
Under the causal-set axiom that information propagates only along
lightlike vectors~\cite{Bombelli1987,Henson2006}, the maximum
number of distinct causal channels available to an observer at the
origin is $\Nloc = |\FS| + |\PS| = 12$.
This count follows from Theorem~\ref{thm:diamond} together with the
stated axiom; it is not a theorem of the present paper.
\end{remark}

\begin{remark}[Kissing Number Coincidence]\label{rem:kissing}
$\Nloc = 12$ coincides numerically with the three-dimensional
Euclidean kissing number $k_3=12$~\cite{SchuetteMaass1953},
but the derivation here uses the Minkowski causal structure of
$\Lat$, not Euclidean sphere packing.
The structural reason for the coincidence is as follows.
The 12 lightlike vectors in $\LL$ are exactly the mixed roots of
$D_4$ (\cref{prop:d4}).
The $D_4$ root system contains the $D_3$ root system as a
sub-lattice; $D_3$ is the face-centred cubic (FCC) lattice, whose
kissing number is 12.
The Minkowski metric $\eta$ selects the $2\times6$ mixed roots of
$D_4$ (one time component, one spatial component each) out of the
full 24-root system, and their count equals $k_3$ by this embedding.
The coincidence does not persist in other spacetime signatures.
\end{remark}

%------------------------------------------------------------------
\section{Lorentzian Boundary Integration}\label{sec:boundary}
%------------------------------------------------------------------

\begin{definition}[Lorentzian Oriented Chain Complex]\label{def:chain}
The oriented chain complex on $(\CD,\BdyCD)$ assigns:
future links $L_f\in\FS$ the orientation $[\bm{0}\to L_f]$
(outgoing);
past links $L_p\in\PS$ the orientation $[L_p\to\bm{0}]$
(incoming).
Past links are incoming because the past null cone
\emph{converges} on the observer at the origin in Lorentzian
geometry, following the standard orientation for causal
boundaries~\cite{Penrose1972, WaldGR}.
Under the Riemannian convention (all links outgoing),
the boundary sum vanishes identically; the Lorentzian
convention is what produces a non-trivial result.
\end{definition}

\begin{theorem}[Lorentzian Boundary Coupling]\label{thm:boundary}
Let $A = A_\mu\,\mathrm{d}x^\mu$ be a $U(1)$ gauge field,
assumed constant at the scale of the causal diamond
(lattice-spacing corrections are of order $(a/\lambda)^2$
for wavelength $\lambda\gg a$).
The discrete boundary sum over $\BdyCD$ with orientations
from Definition~\ref{def:chain} is
\[
  \sum_{L\in\BdyCD} \int_L A
  \;=\; A_\mu\,\neff{}^\mu,
  \qquad \neff = (12,\,0,\,0,\,0).
\]
The spatial components vanish by the six-fold rotational symmetry
of $\FS$; the temporal component equals $12$.
\end{theorem}

\begin{proof}
Under the constant-field assumption, each future link
$[\bm{0}\to L_f]$ contributes $A_\mu L_f^\mu$, and each
past link $[L_p\to\bm{0}]$ contributes $A_\mu(\bm{0}-L_p)^\mu
= -A_\mu L_p^\mu$.
Since $L_p = -L_f$ for each antipodal pair:
\[
  \sum_{L\in\BdyCD}\int_L A
  = \sum_{L_f\in\FS} A_\mu L_f^\mu
  + \sum_{L_p\in\PS}(-A_\mu L_p^\mu)
  = 2\sum_{L_f\in\FS} A_\mu L_f^\mu.
\]
Time component ($\mu=0$): all six $L_f$ have $L_f^0=+1$,
so $2\sum_f L_f^0 = 12$.
Spatial component ($\mu=j$): $\FS$ contains both
$(1,+1,0,0)$ and $(1,-1,0,0)$, and analogously for $j=2,3$,
so $\sum_f L_f^j = 0$.
\qed
\end{proof}

\begin{remark}[Discrete Boundary Sum, Not Stokes Theorem]\label{rem:stokes}
Theorem~\ref{thm:boundary} is a result about the
\emph{discrete oriented boundary sum}
$\sum_{L\in\BdyCD}\int_L A$.
It is not an application of the classical Stokes theorem
$\int_\Omega d\omega = \int_{\partial\Omega}\omega$,
which would require evaluating a bulk curvature integral
$\int_\Omega dA$ over the diamond volume.
No bulk field-strength integral is computed here; the result
is driven entirely by the Lorentzian orientation convention
and the antipodal symmetry of $\BdyCD$.
\end{remark}

\begin{corollary}[Purely Temporal Boundary Source]\label{cor:temporal}
Any action $S_\mathrm{bdy} \propto \sum_{L\in\BdyCD}\int_L A$
derived from the Lorentzian boundary sum generates source terms
only in the time-component equation of motion.
Spatial equations of motion are unaffected.
\end{corollary}

%------------------------------------------------------------------
\section{Plaquette Structure and Algebraic Invariants}\label{sec:plaquettes}
%------------------------------------------------------------------

An \emph{order-$k$ plaquette} is a $k$-element subset of $\LL$
whose vector sum is $\bm{0}$.

\begin{proposition}[Plaquette Enumeration]\label{prop:plaquettes}
\begin{enumerate}[label=(\roman*), noitemsep]
  \item There are no order-3 plaquettes.
  \item There are exactly 21 order-4 plaquettes, all of type
    $2\FS+2\PS$.
  \item There are six order-2 antipodal pairs
    $\{L_f,-L_f\}$ for $L_f\in\FS$.
\end{enumerate}
All 21 quadrilateral plaquettes are listed explicitly in
Appendix~\ref{app:plaquettes}.
\end{proposition}

\begin{proof}
Order-3: any three lightlike vectors have time components
summing to a value in $\{-3,-1,+1,+3\}$, never zero.
Order-4: a zero-sum quadruple requires time components to cancel,
forcing $n_+=n_-=2$ (type $2\FS+2\PS$).
Enumeration of all $\binom{12}{4}=495$ quadruples yields exactly
$21$ satisfying this condition; the full list is in
Appendix~\ref{app:plaquettes}.
\qed
\end{proof}

\begin{remark}[Comparison with CDT]\label{rem:cdt}
Causal dynamical triangulations~\cite{Ambjorn2004} use
triangular simplices (order-3) as fundamental cells.
The absence of order-3 plaquettes in $\CD$ follows directly
from the time-component constraint: a closed triangle on the
causal boundary cannot balance its time components.
The ternary lattice therefore selects order-4 squares as its
minimal closed loops, connecting naturally to Wilson-loop
lattice gauge theory~\cite{Wilson1974} rather than to
simplicial Regge calculus~\cite{Regge1961}.
\end{remark}

\subsection{Incidence Matrix and Plaquette Laplacian}

\begin{definition}[Incidence Matrix and Laplacian]\label{def:laplacian}
Index the twelve links $l=0,\ldots,11$ and the twenty-one
plaquettes $p=0,\ldots,20$.
The \emph{incidence matrix} $\Mmat\in\{0,1\}^{12\times21}$
has $\Mmat_{lp}=1$ iff link $l$ belongs to plaquette $p$.
Each column of $\Mmat$ has exactly four nonzero entries.
The \emph{plaquette Laplacian} is
$\Kmat = \Mmat\Mmat^\top \in\Z^{12\times12}$.
\end{definition}

\begin{proposition}[Plaquette Laplacian Spectrum]\label{prop:spectrum}
The plaquette Laplacian $\Kmat$ has exact integer spectrum
\[
  \mathrm{spec}(\Kmat)
  = \bigl\{0^{(4)},\; 6^{(2)},\; 8^{(3)},\;
            10^{(2)},\; 28^{(1)}\bigr\}.
\]
The four zero eigenvalues span the flat-connection space
$\ker(\Mmat^\top)$ (see \cref{prop:flat});
the eight non-zero eigenvalues describe the curvature degrees
of freedom.
$\rank(\Mmat)=8$; $\Tr(\Kmat)=84=4\times21$.
\end{proposition}

\begin{proof}
Direct diagonalisation of the $12\times12$ integer matrix $\Kmat$
computed from the enumerated incidence matrix.
The explicit matrix structure is described in
Appendix~\ref{app:Kmat}.
All eigenvalues are integers because $\Kmat$ is an integer
symmetric matrix.
\qed
\end{proof}

\subsection{Flat Connections and the Null Space of the Laplacian}
 
\begin{proposition}[Flat-Connection Space]\label{prop:flat}
The plaquette Laplacian $\Kmat = \Mmat\Mmat^\top$ has a
four-dimensional null space:
\[
  \ker(\Kmat) \;=\; \ker(\Mmat^\top)
  \;=\; \bigl\{\theta\in\R^{12} :
  \textstyle\sum_{l\in p}\theta_l = 0
  \;\text{for every plaquette } p\bigr\}.
\]
Elements of $\ker(\Mmat^\top)$ are \emph{flat connections}:
link-variable assignments with vanishing holonomy on every
plaquette.
All four independent flat directions are antipodally antisymmetric:
each satisfies $\theta_{L_f} = -\theta_{-L_f}$ for every antipodal
pair $(L_f, -L_f)$.
The constant assignment $\theta_l = c$ for all $l$ is \emph{not} flat:
it gives holonomy $4$ on every plaquette and is instead an eigenvector
of $\Kmat$ with eigenvalue $28$ (see \cref{prop:quadratic}).
\end{proposition}
 
\begin{proof}
$\ker(\Mmat\Mmat^\top) = \ker(\Mmat^\top)$ because $\Mmat$ has
real entries: $\Mmat\Mmat^\top\theta = 0$ implies
$\theta^\top\Mmat\Mmat^\top\theta = \|\Mmat^\top\theta\|^2 = 0$,
hence $\Mmat^\top\theta = 0$.
The four zero eigenvalues of $\Kmat$
(\cref{prop:spectrum}) give $\dim\ker(\Mmat^\top) = 4$.
Equivalently, $\rank(\Mmat) = 8$ and $\dim\ker(\Mmat^\top)
= 12 - 8 = 4$.
\qed
\end{proof}
 
\begin{remark}[Structure of the Flat-Connection
Space]\label{rem:flat_structure}
The four-dimensional null space $\ker(\Mmat^\top)$ over $\R$
consists of link assignments $\theta$ satisfying the holonomy
constraint $\sum_{l\in p}\theta_l = 0$ simultaneously on all
21 plaquettes.
All four independent flat modes share an antipodal antisymmetry:
each satisfies $\theta_{L_f} = -\theta_{-L_f}$ for every
antipodal pair, so future and past links carry opposite phases.
This structure reflects the Lorentzian time-orientation of
$\BdyCD$: the holonomy constraint forces the future and past
contributions to cancel pairwise.

The constant assignment $\theta_l = c$ does \emph{not} lie in
$\ker(\Mmat^\top)$; it has holonomy $\sum_{l\in p} c = 4c \neq 0$
on every plaquette (for $c\neq0$), confirming it is not a flat
connection.
Gauge fixing (removing the global phase freedom from the path
integral) is therefore a constraint on the integration domain,
not the removal of a zero eigenvalue.

Over $\mathrm{GF}(2)$, the null space $\ker(\Mmat^\top\bmod 2)$
has dimension $5 = 12 - 7$ (since
$\rank(\Mmat)=7$ over $\mathrm{GF}(2)$ versus $8$ over $\R$;
the gap arises from the ternary lattice structure, see
Remark~\ref{rem:CW}).
The $\mathrm{GF}(2)$ null space contains three weight-four vectors
$\{l\in\LL : l_j\neq0\}$ for $j=1,2,3$, corresponding to the
three spatial coordinate axes; these are flat modulo $2$ because
each plaquette intersects each spatial axis in exactly two links,
but they are \emph{not} flat over $\R$ (the holonomy is $2$,
not $0$).
\end{remark}
 
\begin{remark}[CW-Complex Obstruction]\label{rem:CW}
The $1$-skeleton of $\CD$ is a star graph: twelve edges sharing
the single vertex $\bm{0}$.
A star graph is a tree, so its first homology vanishes:
$\ker(\partial_1) = \{0\}$.
The chain-complex condition $\partial_1\circ\partial_2 = 0$
then forces $\mathrm{im}(\partial_2) \subseteq \ker(\partial_1)
= \{0\}$, requiring $\partial_2 = 0$.
But $\rank(\Mmat) = 8 \neq 0$, so the incidence matrix $\Mmat$
cannot serve as a topological boundary operator $\partial_2$ on
this $1$-skeleton.
Concretely, the Euler-characteristic identity
$\chi = \sum(-1)^k\dim C_k = \sum(-1)^k b_k$ would require
$13 - 12 + 21 = 22$ to equal $b_0 - b_1 + b_2$, which is
inconsistent with the v2 claim $(b_0,b_1,b_2) = (1,4,13)$
giving $\chi = 10 \neq 22$.
 
The plaquette Laplacian $\Kmat = \Mmat\Mmat^\top$ and the
flat-connection space $\ker(\Mmat^\top)$ are valid
combinatorial/algebraic objects; their spectra and dimensions
are exact and do not depend on a topological interpretation.
The four-dimensional null space records the gauge-theoretic
flat-connection count, not a Betti number.
\end{remark}
 
\begin{remark}[Extended Boundary Complex]\label{rem:extended}
A valid CW complex can be constructed by replacing the star-graph
$1$-skeleton with the boundary graph whose vertices are the twelve
lightlike vectors and whose edges connect each future--past pair
appearing together in a plaquette.
With the natural alternating cyclic ordering
$f_1\to p_1\to f_2\to p_2$ on each $2\FS+2\PS$ plaquette,
this $1$-skeleton is the complete bipartite graph $K_{6,6}$
(36 edges).
The resulting CW complex has a valid chain complex
$\partial_1\circ\partial_2 = 0$ with cell counts
$|C_0|=12$, $|C_1|=36$, $|C_2|=21$, Betti numbers
$(b_0,b_1,b_2)=(1,7,3)$, and Euler characteristic
$\chi = 1 - 7 + 3 = -3 = 12 - 36 + 21$.
This complex is noted for completeness; the algebraic results
of the present paper (\cref{prop:spectrum,prop:flat}) do not
require it.
\end{remark}

%------------------------------------------------------------------
\section{\texorpdfstring{$U(1)$}{U(1)} Lattice BF Theory}
\label{sec:BF}
%------------------------------------------------------------------

\begin{construction}[$U(1)$ Bulk BF Action]\label{con:BF}
Assign phase $\theta_l\in[0,2\pi)$ to each link.
The holonomy of plaquette $p=(l_1,l_2,l_3,l_4)$ is
$\Theta_p = \theta_{l_1}+\theta_{l_2}+\theta_{l_3}+\theta_{l_4}
\pmod{2\pi}$.
The bulk action is
$S_\mathrm{bulk}(\theta) = -\beta\sum_{p\in\Plaq}\cos(\Theta_p)$.
\end{construction}

\begin{ansatz}[Boundary Chern--Simons Term]\label{ans:bdy}
\ANSATZ\;
Label the six antipodal pairs
$\mathcal{B} = \{(l_f^{(b)}, l_p^{(b)})\}_{b=1}^6$,
with $l_p^{(b)} = -l_f^{(b)}$.
We introduce the boundary action
\[
  S_\mathrm{bdy}(\theta)
  = -\frac{\beta}{\Nloc+1}
    \sum_{b=1}^6
    \cos\!\bigl(\theta_{l_f^{(b)}} - \theta_{l_p^{(b)}}\bigr)
  = -\frac{\beta}{13}
    \sum_{b=1}^6
    \cos\!\bigl(\theta_{l_f^{(b)}} - \theta_{l_p^{(b)}}\bigr).
\]
The sign convention $\theta_{l_f}-\theta_{l_p}$ reflects the
Lorentzian orientation of \cref{thm:boundary}.
The coupling $\beta/13 = \beta/(\Nloc+1)$ is chosen so that
the boundary Laplacian $\Kmat_\mathrm{bdy}$ lifts exactly three
bulk zero modes to eigenvalue $2/13$ (see \cref{prop:quadratic}).
Heuristically, the denominator $\Nloc + 1 = 13$ distributes
the coupling uniformly across the twelve boundary links and the
single internal (observer) degree of freedom at the origin, so
each of the thirteen lattice entities carries equal weight
in the total action.
This choice is not derived from gauge invariance or anomaly
cancellation; deriving the correct coupling from an $\mathrm{SU}(2)$
Plebanski formulation is Open Problem~O2 below.
The full action is $S = S_\mathrm{bulk} + S_\mathrm{bdy}$.
\end{ansatz}

\subsection{Leading-Order Approximate Partition Function}

\begin{proposition}[Character Expansion]\label{prop:character}
Using $e^{\beta\cos\theta} = \sum_{n\in\Z} I_n(\beta)e^{in\theta}$
(where $I_n$ decays rapidly for $|n|>1$ at moderate $\beta$),
the partition function
$\ZBF(\beta) = \int_{[0,2\pi)^{12}}\prod_l \frac{d\theta_l}{2\pi}
e^{-S(\theta)}$
satisfies
\begin{equation}
  \ZBF(\beta) = \sum_{\{n_p\}\in\Z^{21}}
  \prod_{p\in\Plaq} I_{n_p}(\beta)\cdot
  \prod_{l=0}^{11}
  \delta\!\bigl(\textstyle\sum_{p\ni l} n_p,\,0\bigr).
  \label{eq:Zchar}
\end{equation}
The Kronecker delta enforces Gauss's law for each link.
Retaining only $|n_p|\leq1$ gives the leading-order approximation,
justified by the rapid decay of $I_n(\beta)/I_0(\beta)$ for
$|n|>1$.
\end{proposition}

\begin{proposition}[Finite-System Crossover]\label{prop:crossover}
Within the leading-order character truncation of
\cref{prop:character}, the finite-system specific heat
$C(\beta) = \beta^2\,\partial^2(\log\ZBF)/\partial\beta^2$
peaks at $\Cstar\approx 2.7364$ with value $C(\Cstar)\approx 12.69$.
\end{proposition}

\begin{proof}
Truncating to $|n_p|\leq 1$, the zero-mode sector contributes
$Z_0 = I_0(\beta)^{21}$.
For the single-excitation sector, Gauss's law at each link requires
that any excited pair $(n_p=+1, n_q=-1)$ satisfies $\Mmat_{:,p}\cdot\Mmat_{:,q}=0$
(the two plaquettes share no link).
Exhaustive enumeration of all $\binom{21}{2}=210$ plaquette pairs yields
$N_\mathrm{compat}=60$ non-adjacent pairs (verified in Part~7 of the
supplementary code).
Each contributes $I_1(\beta)^2 I_0(\beta)^{19}$, with a factor of two
for the sign choice, giving
\[
  \log \ZBF(\beta)
  \;\approx\;
  21\log I_0(\beta)
  \;+\;
  \log\!\bigl(1 + 120\,\bigl(I_1(\beta)/I_0(\beta)\bigr)^2\bigr).
\]
Numerical differentiation of this expression locates the specific-heat
peak at $\Cstar\approx 2.7364$ with $C(\Cstar)\approx 12.69$.
The computation is reproduced in Part~7 of the supplementary file
\texttt{experiment\_verification\_ID\_paper.py}.
\qed
\end{proof}

\begin{remark}[Finite-Size Caveat]\label{rem:finite}
At $N=12$ links and $21$ plaquettes the configuration space
is finite, so all thermodynamic functions are analytic in $\beta$
and no true phase transition occurs.
The peak at $\Cstar\approx 2.7364$ is a finite-system crossover
with width $\Delta\beta \approx 2/C(\Cstar) \approx 0.158$;
higher Bessel modes ($|n_p|\geq2$) will shift the estimate.
The result is stated for the $|n_p|\leq1$ truncation only.
\end{remark}

\subsection{Saddle-Point Expansion and Spectrum}

\begin{proposition}[Quadratic Action and Spectrum]\label{prop:quadratic}
Expanding around the ordered saddle $\theta_l = 0$ with small
fluctuations $\phi_l$, the quadratic action is
$S_\mathrm{quad} = (\beta/2)\,\bm{\phi}^\top\Kmat_\mathrm{total}
\bm{\phi} + \mathcal{O}(\phi^4)$,
where $\Kmat_\mathrm{total} = \Kmat_\mathrm{bulk} + \Kmat_\mathrm{bdy}$.
 
Labelling links $0,\ldots,5$ as future and $6,\ldots,11$ as their
antipodal past counterparts, the boundary Laplacian is the
$12\times12$ block matrix
\[
  \Kmat_\mathrm{bdy}
  = \frac{1}{13}
  \begin{pmatrix} I_6 & -I_6 \\ -I_6 & I_6 \end{pmatrix},
\]
where $I_6$ is the $6\times6$ identity.
This is the graph Laplacian of the six antipodal pairs, with
coupling $1/13$: expanding $\cos(\theta_{l_f}-\theta_{l_p})
\approx 1 - \tfrac{1}{2}(\theta_{l_f}-\theta_{l_p})^2$ gives
a positive-semidefinite quadratic form whose matrix has
$+1/13$ on the diagonal and $-1/13$ on the off-diagonal entry
connecting each future--past pair.
 
The four zero eigenvalues of $\Kmat_\mathrm{bulk}$ correspond
to antisymmetric flat connections satisfying
$\theta_{L_f} = -\theta_{-L_f}$ for every antipodal pair
(see \cref{rem:flat_structure}).
The boundary Laplacian $\Kmat_\mathrm{bdy}$ acts on these
antisymmetric modes with eigenvalue $2/13$, lifting all four
to that value.
The two eigenvalue-$6$ modes of $\Kmat_\mathrm{bulk}$ are also
antisymmetric and shift by $2/13$;
the eigenvalue-$8$, $10$, and $28$ modes are symmetric
($\theta_{L_f} = +\theta_{-L_f}$) and are annihilated by
$\Kmat_\mathrm{bdy}$.
The exact spectrum of $\Kmat_\mathrm{total}$ is therefore
\[
  \mathrm{spec}(\Kmat_\mathrm{total})
  = \bigl\{(2/13)^{(4)},\; (6+2/13)^{(2)},\;
            8^{(3)},\; 10^{(2)},\; 28^{(1)}\bigr\}.
\]
In particular, $\Kmat_\mathrm{total}$ has no zero eigenvalue:
the global $U(1)$ mode (the all-ones vector) is an eigenvector
of $\Kmat_\mathrm{bulk}$ with eigenvalue $28$
(since each row of $\Kmat$ sums to $28$),
not a zero mode.
It is symmetric and therefore unaffected by $\Kmat_\mathrm{bdy}$.
Gauge fixing removes the constant mode from the path integral,
but this is a constraint on the integration domain, not a
zero-eigenvalue removal.
\end{proposition}

%------------------------------------------------------------------
\section{Discussion}\label{sec:discussion}
%------------------------------------------------------------------

\subsection{Summary}
 
Five exact results have been established.
Lemma~\ref{thm:12lightlike} and Proposition~\ref{prop:d4} are
unconditional mathematical facts, verified by exhaustive finite
enumeration.
Theorem~\ref{thm:diamond} is a verification of five geometric
conditions.
Theorem~\ref{thm:boundary} is a discrete algebra result, valid
under the constant-field assumption stated therein.
Proposition~\ref{prop:spectrum} gives an exact integer eigenvalue
spectrum; Proposition~\ref{prop:flat} identifies the four-dimensional
null space of the plaquette Laplacian with the space of flat
connections on $\CD$.
The star-graph $1$-skeleton of $\CD$ prevents its plaquettes from
being attached as topological $2$-cells, so the null-space dimension
is a gauge-theoretic invariant rather than a Betti number
(Remark~\ref{rem:CW}); an extended CW complex with the correct
homological structure is noted in Remark~\ref{rem:extended}.
Propositions~\ref{prop:character} and~\ref{prop:crossover} give
leading-order approximate results, honestly labelled as such,
with the finite-size caveat of Remark~\ref{rem:finite};
the crossover value $\Cstar\approx 2.7364$ is now reproduced
by explicit computation in Part~7 of the supplementary file.
Ansatz~\ref{ans:bdy} introduces the boundary coupling as a
motivated but underived choice.

\subsection{Open Problems}

\begin{description}[style=nextline]
  \item[(O1) Continuum limit.]
    Does the $U(1)$ BF theory on $\CD$ converge, under Wilsonian
    renormalisation group flow, to a topological field theory in
    the continuum?
    The discrete boundary sum and plaquette spectrum derived here
    are exact at lattice spacing $a = 1$.
    In the limit $\lambda \gg a \to 0$, the finite-system crossover
    at $\Cstar$ must sharpen into a genuine topological phase
    transition for the invariants to acquire continuum significance.
    Establishing this requires embedding the 12-link diamond
    into a sequence of refined lattices and tracking the
    renormalisation flow of both the bulk and boundary couplings.
  \item[(O2) $\mathrm{SU}(2)$ upgrade and Plebanski constraint.]
    Replacing $U(1)$ by $\mathrm{SU}(2)$ and imposing the
    simplicity constraint $B^{IJ}=\ast(e^I\wedge e^J)$ on each
    $2\FS+2\PS$ plaquette may yield a discrete model of
    four-dimensional general relativity.
    The boundary coupling (currently the ansatz $\beta/13$) should
    emerge from this construction; deriving it is the central
    remaining open problem.
  \item[(O3) Topological entanglement entropy.]
    Computing $\gamma=\log D$ from the $\mathrm{SU}(2)$ ground state
    and identifying the Chern--Simons level $k^*$ such that $\gamma=1$
    nat via the Kitaev--Preskill formula~\cite{KitaevPreskill2006}
    are open.
  \item[(O4) Lorentz symmetry restoration.]
    Whether $B_4$ invariance at the lattice scale implies
    $\mathrm{SO}(3,1)$ in the infrared, as observed in causal
    dynamical triangulations~\cite{Ambjorn2004}, is open.
    Dimension-4 Lorentz-violating operators require breaking the
    Minkowski signature (itself $B_4$-invariant); unprotected
    operators first appear at dimension 6, suppressed by $(a/L)^2$.
  \item[(O5) Spectral dimension of the MERA graph.]
    The ternary lattice supports a MERA tensor network with
    coarse-graining base 3.
    The spectral dimension of the resulting causal graph and its
    relationship to conformal factors are left for future work.
\end{description}

%------------------------------------------------------------------
\section*{Acknowledgements}
%------------------------------------------------------------------

The author thanks three anonymous reviewers whose detailed
evaluations prompted the explicit listing of all 21 plaquettes,
the finite-size caveat for the crossover, the distinction between
$\CD$ and $\BdyCD$, the reframing of Theorem~\ref{thm:boundary}
as a discrete boundary sum rather than an application of Stokes'
theorem, and the explicit labelling of the boundary coupling as
an ansatz.
A subsequent review identified the CW-complex obstruction in the
original Betti-number computation, leading to the corrected
flat-connection analysis of v3 (Proposition~\ref{prop:flat} and
Remark~\ref{rem:CW}).

%------------------------------------------------------------------
\begin{thebibliography}{99}

\bibitem{Ambjorn2004}
J.~Ambj\o{}rn, J.~Jurkiewicz and R.~Loll,
Phys.\ Rev.\ Lett.\ \textbf{93} (2004) 131301.

\bibitem{Bombelli1987}
L.~Bombelli et al.,
Phys.\ Rev.\ Lett.\ \textbf{59} (1987) 521.

\bibitem{Henson2006}
J.~Henson, in \emph{Approaches to Quantum Gravity},
Cambridge University Press (2009).

\bibitem{KitaevPreskill2006}
A.~Kitaev and J.~Preskill,
Phys.\ Rev.\ Lett.\ \textbf{96} (2006) 110404.

\bibitem{Kogut1979}
J.~B.~Kogut, Rev.\ Mod.\ Phys.\ \textbf{51} (1979) 659.

\bibitem{Penrose1972}
R.~Penrose, \emph{Techniques of Differential Topology in
Relativity}, SIAM (1972).

\bibitem{Perez2013}
A.~Perez, Living Rev.\ Relat.\ \textbf{16} (2013) 3.

\bibitem{Regge1961}
T.~Regge, Nuovo Cimento \textbf{19} (1961) 558.

\bibitem{SchuetteMaass1953}
C.~Sch{\"u}tte and B.~L.~van der Waerden,
Math.\ Ann.\ \textbf{125} (1953) 325.

\bibitem{WaldGR}
R.~M.~Wald, \emph{General Relativity},
University of Chicago Press (1984).

\bibitem{Wilson1974}
K.~G.~Wilson, Phys.\ Rev.\ D \textbf{10} (1974) 2445.

\end{thebibliography}

%------------------------------------------------------------------
\appendix
%------------------------------------------------------------------

\section{Explicit Enumeration of Lightlike Vectors}\label{app:enum}

\noindent\textbf{Future null sheet} $\FS$ ($v_0=+1$):
\[
  (1,{-}1,0,0),\;(1,{+}1,0,0),\;(1,0,{-}1,0),\;
  (1,0,{+}1,0),\;(1,0,0,{-}1),\;(1,0,0,{+}1).
\]
\textbf{Past null sheet} $\PS$ ($v_0=-1$):
\[
  ({-}1,{-}1,0,0),\;({-}1,{+}1,0,0),\;({-}1,0,{-}1,0),\;
  ({-}1,0,{+}1,0),\;({-}1,0,0,{-}1),\;({-}1,0,0,{+}1).
\]

\section{Pairwise Spacelike Verification for $\FS$}\label{app:pairwise}

For distinct $u,v\in\FS$: $(u-v)_0=0$ and at least one spatial
component differs, so $\ssp(u-v)\geq2>0$. All 15 pairs verified.

\section{Spanning of $\R^4$ by $\FS$}\label{app:span}

Four vectors $w_1=(1,+1,0,0)$, $w_2=(1,-1,0,0)$,
$w_3=(1,0,+1,0)$, $w_4=(1,0,0,+1)$ have $\det W = -2\neq0$;
hence $\rank(W)=4$ and $\FS$ spans $\R^4$.

\section{All 21 Order-4 Plaquettes}\label{app:plaquettes}

Each row gives four lightlike vectors summing to $\bm{0}$.
All plaquettes are of type $2\FS+2\PS$
(two future, two past links).

\begin{longtable}{cl}
\toprule
\# & Links $(v_1, v_2, v_3, v_4)$\\
\midrule
\endhead
\bottomrule
\endfoot
1  & $(1,-1,0,0),\;(1,+1,0,0),\;(-1,-1,0,0),\;(-1,+1,0,0)$\\
2  & $(1,-1,0,0),\;(1,+1,0,0),\;(-1,0,-1,0),\;(-1,0,+1,0)$\\
3  & $(1,-1,0,0),\;(1,+1,0,0),\;(-1,0,0,-1),\;(-1,0,0,+1)$\\
4  & $(1,-1,0,0),\;(1,0,-1,0),\;(-1,+1,0,0),\;(-1,0,+1,0)$\\
5  & $(1,-1,0,0),\;(1,0,+1,0),\;(-1,+1,0,0),\;(-1,0,-1,0)$\\
6  & $(1,-1,0,0),\;(1,0,0,-1),\;(-1,+1,0,0),\;(-1,0,0,+1)$\\
7  & $(1,-1,0,0),\;(1,0,0,+1),\;(-1,+1,0,0),\;(-1,0,0,-1)$\\
8  & $(1,+1,0,0),\;(1,0,-1,0),\;(-1,-1,0,0),\;(-1,0,+1,0)$\\
9  & $(1,+1,0,0),\;(1,0,+1,0),\;(-1,-1,0,0),\;(-1,0,-1,0)$\\
10 & $(1,+1,0,0),\;(1,0,0,-1),\;(-1,-1,0,0),\;(-1,0,0,+1)$\\
11 & $(1,+1,0,0),\;(1,0,0,+1),\;(-1,-1,0,0),\;(-1,0,0,-1)$\\
12 & $(1,0,-1,0),\;(1,0,+1,0),\;(-1,-1,0,0),\;(-1,+1,0,0)$\\
13 & $(1,0,-1,0),\;(1,0,+1,0),\;(-1,0,0,-1),\;(-1,0,0,+1)$\\
14 & $(1,0,-1,0),\;(1,0,0,-1),\;(-1,0,+1,0),\;(-1,0,0,+1)$\\
15 & $(1,0,-1,0),\;(1,0,0,+1),\;(-1,0,+1,0),\;(-1,0,0,-1)$\\
16 & $(1,0,+1,0),\;(1,0,0,-1),\;(-1,0,-1,0),\;(-1,0,0,+1)$\\
17 & $(1,0,+1,0),\;(1,0,0,+1),\;(-1,0,-1,0),\;(-1,0,0,-1)$\\
18 & $(1,0,0,-1),\;(1,0,0,+1),\;(-1,-1,0,0),\;(-1,+1,0,0)$\\
19 & $(1,0,0,-1),\;(1,0,0,+1),\;(-1,0,-1,0),\;(-1,0,+1,0)$\\
20 & $(1,0,-1,0),\;(1,0,+1,0),\;(-1,0,-1,0),\;(-1,0,+1,0)$\\
21 & $(1,0,0,-1),\;(1,0,0,+1),\;(-1,0,0,-1),\;(-1,0,0,+1)$\\
\end{longtable}

\section{Structure of the Plaquette Laplacian}\label{app:Kmat}

By the $2\FS+2\PS$ structure, every link participates in the
same number of plaquettes, giving $\Kmat_{ll}=7$ for all $l$
(since $\Tr(\Kmat)=84$ and there are 12 links).
Off-diagonal entries $\Kmat_{ll'}$ count shared plaquettes,
taking values in $\{0,1,2,3\}$.
The full $12\times12$ integer matrix and the code used for its
diagonalisation are available in the supplementary file
\texttt{experiment\_verification\_ID\_paper.py} via the command
\texttt{K = M @ M.T} where \texttt{M} is the $12\times21$
incidence matrix computed from the plaquettes of
Appendix~\ref{app:plaquettes}.
The character-expansion computation of \cref{prop:crossover}
(60 compatible plaquette pairs, specific-heat peak at
$\Cstar\approx 2.7364$) is reproduced in Part~7 of the same file.

\end{document}